\subsection{p+ Implant}
\begin{center}
	\begin{tikzpicture}[node distance = 3cm, auto, thick,scale=\CrossAndTopSectionBig, every node/.style={transform shape}]
		% substrate
\fill[YellowOrange] (0,0) rectangle (20,2);
% n-well
\fill[Goldenrod] (1.25,0.75) rectangle (8.25,2);
% source
\fill[RedOrange] (0,1.5) rectangle (1,2);
\fill[RedOrange] (3.5,1) rectangle (5,2);
\fill[RedOrange] (6.5,1) rectangle (8,2);
\fill[RedOrange] (8.5,1.5) rectangle (11.5,2);
\fill[RedOrange] (17,1) rectangle (18.5,2);
\fill[RedOrange] (19,1.5) rectangle (20,2);
% gate oxide
\fill[LightGray] (4.8,2) rectangle (6.7,2.1);
\fill[LightGray] (13.3,2) rectangle (15.2,2.1);
% gate poly
\fill[BrickRed] (4.8,2.1) rectangle (6.7,2.2);
\fill[BrickRed] (13.3,2.1) rectangle (15.2,2.2);
%field oxides:
\fill[DarkGray] (0,2) rectangle (1,4);
\fill[DarkGray] (8.5,2) rectangle (11.5,4);
\fill[DarkGray] (19,2) rectangle (20,4);
% channel stop
\fill[RedOrange] (0,1.5) rectangle (1,2);
\fill[RedOrange] (8.5,1.5) rectangle (11.5,2);
\fill[RedOrange] (19,1.5) rectangle (20,2);
	\end{tikzpicture}
	\begin{tikzpicture}[node distance = 3cm, auto, thick,scale=\CrossAndTopSectionBig, every node/.style={transform shape}]
		\fill[YellowOrange] (0,0) rectangle (20,10);
% n-well
\fill[Goldenrod] (1.25,1.5) rectangle (8.25,7.25);
% p+
\fill[RedOrange] (3.5,2) rectangle (5,6.5);
\fill[RedOrange] (6.5,2) rectangle (8,6.5);
\fill[RedOrange] (17,2) rectangle (18.5,6.5);
% field oxide
\fill[DarkGray] (0,0) rectangle (1,12);
\fill[DarkGray] (8.5,0) rectangle (11.5,12);
\fill[DarkGray] (19,0) rectangle (20,12);
\fill[DarkGray] (0,0) rectangle (20,1.25);
\fill[DarkGray] (0,7.5) rectangle (20,12);
% poly
\fill[BrickRed] (4.8,1.75) rectangle (6.7,9);
\fill[BrickRed] (13.3,1.75) rectangle (15.2,9);
\fill[BrickRed] (4.8,8) rectangle (15.2,9);
	\end{tikzpicture}
\end{center}

For the bulk of the NMOS transistors and for the source and drain of the PMOS transistors highly doped  p+ areas are required.
In this step we're going to build these.

\subsubsection{Dioxide layer}
\begin{center}
	\begin{tikzpicture}[node distance = 3cm, auto, thick,scale=\CrossSectionOnly, every node/.style={transform shape}]
		% substrate
\fill[YellowOrange] (0,0) rectangle (20,2);
\node at (2,0.5) {Si (p-type)};
% n-well
\fill[Goldenrod] (1.25,0.75) rectangle (8.25,2);
\node at (4.75,1) {N-Well};
% gate oxide
\fill[LightGray] (4.8,2) rectangle (6.7,2.1);
\fill[LightGray] (13.3,2) rectangle (15.2,2.1);
% gate poly
\fill[BrickRed] (4.8,2.1) rectangle (6.7,2.2);
\fill[BrickRed] (13.3,2.1) rectangle (15.2,2.2);
%field oxides:
\fill[DarkGray] (0,2) rectangle (1,4);
\fill[DarkGray] (8.5,2) rectangle (11.5,4);
\fill[DarkGray] (19,2) rectangle (20,4);
% channel stop
\fill[RedOrange] (0,1.5) rectangle (1,2);
\fill[RedOrange] (8.5,1.5) rectangle (11.5,2);
\fill[RedOrange] (19,1.5) rectangle (20,2);
	\end{tikzpicture} \\
	\includegraphics[scale=0.01]{down_arrow.png} \\
	\begin{tikzpicture}[node distance = 3cm, auto, thick,scale=\CrossSectionOnly, every node/.style={transform shape}]
		% substrate
\fill[YellowOrange] (0,0) rectangle (20,2);
\node at (2,0.5) {Si (p-type)};
% n-well
\fill[Goldenrod] (1.25,0.75) rectangle (8.25,2);
\node at (4.75,1) {N-Well};
% gate oxide
\fill[LightGray] (4.8,2) rectangle (6.7,2.1);
\fill[LightGray] (13.3,2) rectangle (15.2,2.1);
% gate poly
\fill[BrickRed] (4.8,2.1) rectangle (6.7,2.2);
\fill[BrickRed] (13.3,2.1) rectangle (15.2,2.2);
% dioxide
\fill[NormalGray] (0,2) rectangle (4.8,2.7);
\fill[NormalGray] (4.8,2.2) rectangle (6.7,2.7);
\fill[NormalGray] (6.7,2) rectangle (13.3,2.7);
\fill[NormalGray] (13.3,2.2) rectangle (15.2,2.7);
\fill[NormalGray] (15.2,2) rectangle (20,2.7);
%field oxides:
\fill[DarkGray] (0,2) rectangle (1,4);
\fill[DarkGray] (8.5,2) rectangle (11.5,4);
\fill[DarkGray] (19,2) rectangle (20,4);
% channel stop
\fill[RedOrange] (0,1.5) rectangle (1,2);
\fill[RedOrange] (8.5,1.5) rectangle (11.5,2);
\fill[RedOrange] (19,1.5) rectangle (20,2);
	\end{tikzpicture}
\end{center}

\subsubsection{Pattering}
\begin{center}
	\begin{tikzpicture}[node distance = 3cm, auto, thick,scale=\CrossAndTopSection, every node/.style={transform shape}]
		% substrate
\fill[YellowOrange] (0,0) rectangle (20,2);
\node at (2,0.5) {Si (p-type)};
% n-well
\fill[Goldenrod] (1.25,0.75) rectangle (8.25,2);
\node at (4.75,1) {N-Well};
% gate oxide
\fill[LightGray] (4.8,2) rectangle (6.7,2.1);
\fill[LightGray] (13.3,2) rectangle (15.2,2.1);
% gate poly
\fill[BrickRed] (4.8,2.1) rectangle (6.7,2.2);
\fill[BrickRed] (13.3,2.1) rectangle (15.2,2.2);
% dioxide
\fill[NormalGray] (0,2) rectangle (4.8,2.7);
\fill[NormalGray] (4.8,2.2) rectangle (6.7,2.7);
\fill[NormalGray] (6.7,2) rectangle (13.3,2.7);
\fill[NormalGray] (13.3,2.2) rectangle (15.2,2.7);
\fill[NormalGray] (15.2,2) rectangle (20,2.7);
%field oxides:
\fill[DarkGray] (0,2) rectangle (1,4);
\fill[DarkGray] (8.5,2) rectangle (11.5,4);
\fill[DarkGray] (19,2) rectangle (20,4);
% channel stop
\fill[RedOrange] (0,1.5) rectangle (1,2);
\fill[RedOrange] (8.5,1.5) rectangle (11.5,2);
\fill[RedOrange] (19,1.5) rectangle (20,2);
	\end{tikzpicture}
	\begin{tikzpicture}[node distance = 3cm, auto, thick,scale=\CrossAndTopSection, every node/.style={transform shape}]
		\fill[NormalGray] (0,0) rectangle (20,12);
	\end{tikzpicture} \\
	\includegraphics[scale=0.01]{down_arrow.png} \\
	\begin{tikzpicture}[node distance = 3cm, auto, thick,scale=\CrossAndTopSection, every node/.style={transform shape}]
		% substrate
\fill[YellowOrange] (0,0) rectangle (20,2);
\node at (2,0.5) {Si (p-type)};
% n-well
\fill[Goldenrod] (1.25,0.75) rectangle (8.25,2);
\node at (4.75,1) {N-Well};
% gate oxide
\fill[LightGray] (4.8,2) rectangle (6.7,2.1);
\fill[LightGray] (13.3,2) rectangle (15.2,2.1);
% gate poly
\fill[BrickRed] (4.8,2.1) rectangle (6.7,2.2);
\fill[BrickRed] (13.3,2.1) rectangle (15.2,2.2);
% dioxide
\fill[NormalGray] (0,2) rectangle (4.8,2.7);
\fill[NormalGray] (4.8,2.2) rectangle (6.7,2.7);
\fill[NormalGray] (6.7,2) rectangle (13.3,2.7);
\fill[NormalGray] (13.3,2.2) rectangle (15.2,2.7);
\fill[NormalGray] (15.2,2) rectangle (20,2.7);
% resist
\fill[orange] (0,4) rectangle (1,4.6);
\fill[orange] (1,2.7) rectangle (3.5,4.6);
\fill[orange] (8,2.7) rectangle (8.5,4.6);
\fill[orange] (8.5,4) rectangle (11.5,4.6);
\fill[orange] (11.5,2.7) rectangle (17,4.6);
\fill[orange] (18.5,2.7) rectangle (19,4.6);
\fill[orange] (19,4) rectangle (20,4.6);
%field oxides:
\fill[DarkGray] (0,2) rectangle (1,4);
\fill[DarkGray] (8.5,2) rectangle (11.5,4);
\fill[DarkGray] (19,2) rectangle (20,4);
% channel stop
\fill[RedOrange] (0,1.5) rectangle (1,2);
\fill[RedOrange] (8.5,1.5) rectangle (11.5,2);
\fill[RedOrange] (19,1.5) rectangle (20,2);
	\end{tikzpicture}
	\begin{tikzpicture}[node distance = 3cm, auto, thick,scale=\CrossAndTopSection, every node/.style={transform shape}]
		\fill[orange] (0,0) rectangle (20,12);
\fill[NormalGray] (3.5,2) rectangle (8,6.5);
\fill[NormalGray] (17,2) rectangle (18.5,6.5);
	\end{tikzpicture}
\end{center}

\subsubsection{Etching}
\begin{center}
	\begin{tikzpicture}[node distance = 3cm, auto, thick,scale=\CrossAndTopSection, every node/.style={transform shape}]
		% substrate
\fill[YellowOrange] (0,0) rectangle (20,2);
\node at (2,0.5) {Si (p-type)};
% n-well
\fill[Goldenrod] (1.25,0.75) rectangle (8.25,2);
\node at (4.75,1) {N-Well};
% gate oxide
\fill[LightGray] (4.8,2) rectangle (6.7,2.1);
\fill[LightGray] (13.3,2) rectangle (15.2,2.1);
% gate poly
\fill[BrickRed] (4.8,2.1) rectangle (6.7,2.2);
\fill[BrickRed] (13.3,2.1) rectangle (15.2,2.2);
% dioxide
\fill[NormalGray] (0,2) rectangle (4.8,2.7);
\fill[NormalGray] (4.8,2.2) rectangle (6.7,2.7);
\fill[NormalGray] (6.7,2) rectangle (13.3,2.7);
\fill[NormalGray] (13.3,2.2) rectangle (15.2,2.7);
\fill[NormalGray] (15.2,2) rectangle (20,2.7);
% resist
\fill[orange] (0,4) rectangle (1,4.6);
\fill[orange] (1,2.7) rectangle (3.5,4.6);
\fill[orange] (8,2.7) rectangle (8.5,4.6);
\fill[orange] (8.5,4) rectangle (11.5,4.6);
\fill[orange] (11.5,2.7) rectangle (17,4.6);
\fill[orange] (18.5,2.7) rectangle (19,4.6);
\fill[orange] (19,4) rectangle (20,4.6);
%field oxides:
\fill[DarkGray] (0,2) rectangle (1,4);
\fill[DarkGray] (8.5,2) rectangle (11.5,4);
\fill[DarkGray] (19,2) rectangle (20,4);
% channel stop
\fill[RedOrange] (0,1.5) rectangle (1,2);
\fill[RedOrange] (8.5,1.5) rectangle (11.5,2);
\fill[RedOrange] (19,1.5) rectangle (20,2);
	\end{tikzpicture}
	\begin{tikzpicture}[node distance = 3cm, auto, thick,scale=\CrossAndTopSection, every node/.style={transform shape}]
		\fill[orange] (0,0) rectangle (20,12);
\fill[NormalGray] (3.5,2) rectangle (8,6.5);
\fill[NormalGray] (17,2) rectangle (18.5,6.5);
	\end{tikzpicture} \\
	\includegraphics[scale=0.01]{down_arrow.png} \\
	\begin{tikzpicture}[node distance = 3cm, auto, thick,scale=\CrossAndTopSection, every node/.style={transform shape}]
		% substrate
\fill[YellowOrange] (0,0) rectangle (20,2);
\node at (2,0.5) {Si (p-type)};
% n-well
\fill[Goldenrod] (1.25,0.75) rectangle (8.25,2);
\node at (4.75,1) {N-Well};
% gate oxide
\fill[LightGray] (4.8,2) rectangle (6.7,2.1);
\fill[LightGray] (13.3,2) rectangle (15.2,2.1);
% gate poly
\fill[BrickRed] (4.8,2.1) rectangle (6.7,2.2);
\fill[BrickRed] (13.3,2.1) rectangle (15.2,2.2);
% dioxide
\fill[NormalGray] (0,2) rectangle (3.5,2.7);
\fill[NormalGray] (8,2) rectangle (13.3,2.7);
\fill[NormalGray] (13.3,2.2) rectangle (15.2,2.7);
\fill[NormalGray] (15.2,2) rectangle (17,2.7);
\fill[NormalGray] (18.5,2) rectangle (20,2.7);
% resist
\fill[orange] (0,4) rectangle (1,4.6);
\fill[orange] (1,2.7) rectangle (3.5,3.3);
\fill[orange] (8,2.7) rectangle (8.5,3.3);
\fill[orange] (8.5,4) rectangle (11.5,4.6);
\fill[orange] (11.5,2.7) rectangle (17,3.3);
\fill[orange] (18.5,2.7) rectangle (19,3.3);
\fill[orange] (19,4) rectangle (20,4.6);
%field oxides:
\fill[DarkGray] (0,2) rectangle (1,4);
\fill[DarkGray] (8.5,2) rectangle (11.5,4);
\fill[DarkGray] (19,2) rectangle (20,4);
% channel stop
\fill[RedOrange] (0,1.5) rectangle (1,2);
\fill[RedOrange] (8.5,1.5) rectangle (11.5,2);
\fill[RedOrange] (19,1.5) rectangle (20,2);
	\end{tikzpicture}
	\begin{tikzpicture}[node distance = 3cm, auto, thick,scale=\CrossAndTopSection, every node/.style={transform shape}]
		\fill[orange] (0,0) rectangle (20,12);
% source & drain
\fill[Goldenrod] (2.5,1.5) rectangle (7,6);
% body
\fill[YellowOrange] (17,1.5) rectangle (18.5,6);
% gate poly
\fill[BrickRed] (3.8,1.5) rectangle (5.7,6);
	\end{tikzpicture}
\end{center}

\subsubsection{Cleaning}
\begin{center}
	\begin{tikzpicture}[node distance = 3cm, auto, thick,scale=\CrossSectionOnly, every node/.style={transform shape}]
		% substrate
\fill[YellowOrange] (0,0) rectangle (20,2);
\node at (2,0.5) {Si (p-type)};
% n-well
\fill[Goldenrod] (1.25,0.75) rectangle (8.25,2);
\node at (4.75,1) {N-Well};
% gate oxide
\fill[LightGray] (4.8,2) rectangle (6.7,2.1);
\fill[LightGray] (13.3,2) rectangle (15.2,2.1);
% gate poly
\fill[BrickRed] (4.8,2.1) rectangle (6.7,2.2);
\fill[BrickRed] (13.3,2.1) rectangle (15.2,2.2);
% dioxide
\fill[NormalGray] (0,2) rectangle (3.5,2.7);
\fill[NormalGray] (8,2) rectangle (13.3,2.7);
\fill[NormalGray] (13.3,2.2) rectangle (15.2,2.7);
\fill[NormalGray] (15.2,2) rectangle (17,2.7);
\fill[NormalGray] (18.5,2) rectangle (20,2.7);
% resist
\fill[orange] (0,4) rectangle (1,4.6);
\fill[orange] (1,2.7) rectangle (3.5,4.6);
\fill[orange] (8,2.7) rectangle (8.5,4.6);
\fill[orange] (8.5,4) rectangle (11.5,4.6);
\fill[orange] (11.5,2.7) rectangle (17,4.6);
\fill[orange] (18.5,2.7) rectangle (19,4.6);
\fill[orange] (19,4) rectangle (20,4.6);
%field oxides:
\fill[DarkGray] (0,2) rectangle (1,4);
\fill[DarkGray] (8.5,2) rectangle (11.5,4);
\fill[DarkGray] (19,2) rectangle (20,4);
% channel stop
\fill[RedOrange] (0,1.5) rectangle (1,2);
\fill[RedOrange] (8.5,1.5) rectangle (11.5,2);
\fill[RedOrange] (19,1.5) rectangle (20,2);
	\end{tikzpicture} \\
	\includegraphics[scale=0.01]{down_arrow.png} \\
	\begin{tikzpicture}[node distance = 3cm, auto, thick,scale=\CrossSectionOnly, every node/.style={transform shape}]
		% substrate
\fill[YellowOrange] (0,0) rectangle (20,2);
\node at (2,0.5) {Si (p-type)};
% n-well
\fill[Goldenrod] (1.25,0.75) rectangle (8.25,2);
\node at (4.75,1) {N-Well};
% gate oxide
\fill[LightGray] (4.8,2) rectangle (6.7,2.1);
\fill[LightGray] (13.3,2) rectangle (15.2,2.1);
% gate poly
\fill[BrickRed] (4.8,2.1) rectangle (6.7,2.2);
\fill[BrickRed] (13.3,2.1) rectangle (15.2,2.2);
% dioxide
\fill[NormalGray] (0,2) rectangle (3.5,2.7);
\fill[NormalGray] (8,2) rectangle (13.3,2.7);
\fill[NormalGray] (13.3,2.2) rectangle (15.2,2.7);
\fill[NormalGray] (15.2,2) rectangle (17,2.7);
\fill[NormalGray] (18.5,2) rectangle (20,2.7);
%field oxides:
\fill[DarkGray] (0,2) rectangle (1,4);
\fill[DarkGray] (8.5,2) rectangle (11.5,4);
\fill[DarkGray] (19,2) rectangle (20,4);
% channel stop
\fill[RedOrange] (0,1.5) rectangle (1,2);
\fill[RedOrange] (8.5,1.5) rectangle (11.5,2);
\fill[RedOrange] (19,1.5) rectangle (20,2);
	\end{tikzpicture}
\end{center}

\subsubsection{Predeposition}
\begin{center}
	\begin{tikzpicture}[node distance = 3cm, auto, thick,scale=\CrossSectionOnly, every node/.style={transform shape}]
		% substrate
\fill[YellowOrange] (0,0) rectangle (20,2);
\node at (2,0.5) {Si (p-type)};
% n-well
\fill[Goldenrod] (1.25,0.75) rectangle (8.25,2);
\node at (4.75,1) {N-Well};
% gate oxide
\fill[LightGray] (4.8,2) rectangle (6.7,2.1);
\fill[LightGray] (13.3,2) rectangle (15.2,2.1);
% gate poly
\fill[BrickRed] (4.8,2.1) rectangle (6.7,2.2);
\fill[BrickRed] (13.3,2.1) rectangle (15.2,2.2);
% dioxide
\fill[NormalGray] (0,2) rectangle (3.5,2.7);
\fill[NormalGray] (8,2) rectangle (13.3,2.7);
\fill[NormalGray] (13.3,2.2) rectangle (15.2,2.7);
\fill[NormalGray] (15.2,2) rectangle (17,2.7);
\fill[NormalGray] (18.5,2) rectangle (20,2.7);
%field oxides:
\fill[DarkGray] (0,2) rectangle (1,4);
\fill[DarkGray] (8.5,2) rectangle (11.5,4);
\fill[DarkGray] (19,2) rectangle (20,4);
% channel stop
\fill[RedOrange] (0,1.5) rectangle (1,2);
\fill[RedOrange] (8.5,1.5) rectangle (11.5,2);
\fill[RedOrange] (19,1.5) rectangle (20,2);

\forloop{ct}{0}{\value{ct} < 21}
{
	\draw [->] (\value{ct},5) -- (\value{ct},4.2);
	\node at (\value{ct},5.2) {B$^{11}$};
}
	\end{tikzpicture} \\
	\includegraphics[scale=0.01]{down_arrow.png} \\
	\begin{tikzpicture}[node distance = 3cm, auto, thick,scale=\CrossSectionOnly, every node/.style={transform shape}]
		% substrate
\fill[YellowOrange] (0,0) rectangle (20,2);
\node at (2,0.5) {Si (p-type)};
% n-well
\fill[Goldenrod] (1.25,0.75) rectangle (8.25,2);
\node at (4.75,1) {N-Well};
% gate oxide
\fill[LightGray] (4.8,2) rectangle (6.7,2.1);
\fill[LightGray] (13.3,2) rectangle (15.2,2.1);
% gate poly
\fill[BrickRed] (4.8,2.1) rectangle (6.7,2.2);
\fill[BrickRed] (13.3,2.1) rectangle (15.2,2.2);
% dioxide
\fill[NormalGray] (0,2) rectangle (3.5,2.7);
\fill[NormalGray] (8,2) rectangle (13.3,2.7);
\fill[NormalGray] (13.3,2.2) rectangle (15.2,2.7);
\fill[NormalGray] (15.2,2) rectangle (17,2.7);
\fill[NormalGray] (18.5,2) rectangle (20,2.7);
% p+
\fill[RedOrange] (0,4) rectangle (1,4.1);
\fill[RedOrange] (1,2.7) rectangle (3.5,2.8);
\fill[RedOrange] (8,2.7) rectangle (8.5,2.8);
\fill[RedOrange] (8.5,4) rectangle (11.5,4.1);
\fill[RedOrange] (11.5,2.7) rectangle (17,2.8);
\fill[RedOrange] (18.5,2.7) rectangle (19,2.8);
\fill[RedOrange] (19,4) rectangle (20,4.1);
\fill[RedOrange] (3.5,1.8) rectangle (4.8,2);
\fill[RedOrange] (4.8,2.2) rectangle (6.7,2.3);
\fill[RedOrange] (6.7,1.8) rectangle (8,2);
\fill[RedOrange] (17,1.8) rectangle (18.5,2);
%field oxides:
\fill[DarkGray] (0,2) rectangle (1,4);
\fill[DarkGray] (8.5,2) rectangle (11.5,4);
\fill[DarkGray] (19,2) rectangle (20,4);
% channel stop
\fill[RedOrange] (0,1.5) rectangle (1,2);
\fill[RedOrange] (8.5,1.5) rectangle (11.5,2);
\fill[RedOrange] (19,1.5) rectangle (20,2);
	\end{tikzpicture}
\end{center}

The p+ islands are implanted with a Boron ($B^{11}$) dose of $4\times10^{11}cm^{-2}$ at an energy of 35 KeV.

\subsubsection{Sacrificial oxide}
\begin{center}
	\begin{tikzpicture}[node distance = 3cm, auto, thick,scale=\CrossSectionOnly, every node/.style={transform shape}]
		% substrate
\fill[YellowOrange] (0,0) rectangle (20,2);
\node at (2,0.5) {Si (p-type)};
% n-well
\fill[Goldenrod] (1.25,0.75) rectangle (8.25,2);
\node at (4.75,1) {N-Well};
% gate oxide
\fill[LightGray] (4.8,2) rectangle (6.7,2.1);
\fill[LightGray] (13.3,2) rectangle (15.2,2.1);
% gate poly
\fill[BrickRed] (4.8,2.1) rectangle (6.7,2.2);
\fill[BrickRed] (13.3,2.1) rectangle (15.2,2.2);
% dioxide
\fill[NormalGray] (0,2) rectangle (3.5,2.7);
\fill[NormalGray] (8,2) rectangle (13.3,2.7);
\fill[NormalGray] (13.3,2.2) rectangle (15.2,2.7);
\fill[NormalGray] (15.2,2) rectangle (17,2.7);
\fill[NormalGray] (18.5,2) rectangle (20,2.7);
% p+
\fill[RedOrange] (0,4) rectangle (1,4.1);
\fill[RedOrange] (1,2.7) rectangle (3.5,2.8);
\fill[RedOrange] (8,2.7) rectangle (8.5,2.8);
\fill[RedOrange] (8.5,4) rectangle (11.5,4.1);
\fill[RedOrange] (11.5,2.7) rectangle (17,2.8);
\fill[RedOrange] (18.5,2.7) rectangle (19,2.8);
\fill[RedOrange] (19,4) rectangle (20,4.1);
\fill[RedOrange] (3.5,1.8) rectangle (4.8,2);
\fill[RedOrange] (4.8,2.2) rectangle (6.7,2.3);
\fill[RedOrange] (6.7,1.8) rectangle (8,2);
\fill[RedOrange] (17,1.8) rectangle (18.5,2);
%field oxides:
\fill[DarkGray] (0,2) rectangle (1,4);
\fill[DarkGray] (8.5,2) rectangle (11.5,4);
\fill[DarkGray] (19,2) rectangle (20,4);
% channel stop
\fill[RedOrange] (0,1.5) rectangle (1,2);
\fill[RedOrange] (8.5,1.5) rectangle (11.5,2);
\fill[RedOrange] (19,1.5) rectangle (20,2);
	\end{tikzpicture} \\
	\includegraphics[scale=0.01]{down_arrow.png} \\
	\begin{tikzpicture}[node distance = 3cm, auto, thick,scale=\CrossSectionOnly, every node/.style={transform shape}]
		% substrate
\fill[YellowOrange] (0,0) rectangle (20,2);
\node at (2,0.5) {Si (p-type)};
% n-well
\fill[Goldenrod] (1.25,0.75) rectangle (8.25,2);
\node at (4.75,1) {N-Well};
% gate oxide
\fill[LightGray] (4.8,2) rectangle (6.7,2.1);
\fill[LightGray] (13.3,2) rectangle (15.2,2.1);
% gate poly
\fill[BrickRed] (4.8,2.1) rectangle (6.7,2.2);
\fill[BrickRed] (13.3,2.1) rectangle (15.2,2.2);
% dioxide
\fill[NormalGray] (0,2) rectangle (3.5,2.7);
\fill[NormalGray] (8,2) rectangle (13.3,2.7);
\fill[NormalGray] (13.3,2.2) rectangle (15.2,2.7);
\fill[NormalGray] (15.2,2) rectangle (17,2.7);
\fill[NormalGray] (18.5,2) rectangle (20,2.7);
% dioxide 2nd
\fill[NormalGray] (0,2.8) rectangle (3.5,3.3);
\fill[NormalGray] (3.5,2) rectangle (4.8,2.8);
\fill[NormalGray] (4.8,2.3) rectangle (6.7,2.8);
\fill[NormalGray] (6.7,2) rectangle (8,2.8);
\fill[NormalGray] (8,2.8) rectangle (8.5,3.3);
\fill[NormalGray] (11.5,2.8) rectangle (17,3.3);
\fill[NormalGray] (17,2) rectangle (18.5,2.8);
\fill[NormalGray] (18.5,2.8) rectangle (19,3.3);

\fill[DarkGray] (0,4.1) rectangle (1,4.2);
\fill[DarkGray] (8.5,4.1) rectangle (11.5,4.2);
\fill[DarkGray] (19,4.1) rectangle (20,4.2);
% p+
\fill[RedOrange] (0,4) rectangle (1,4.1);
\fill[RedOrange] (1,2.7) rectangle (3.5,2.8);
\fill[RedOrange] (8,2.7) rectangle (8.5,2.8);
\fill[RedOrange] (8.5,4) rectangle (11.5,4.1);
\fill[RedOrange] (11.5,2.7) rectangle (17,2.8);
\fill[RedOrange] (18.5,2.7) rectangle (19,2.8);
\fill[RedOrange] (19,4) rectangle (20,4.1);
\fill[RedOrange] (3.5,1.8) rectangle (4.8,2);
\fill[RedOrange] (4.8,2.2) rectangle (6.7,2.3);
\fill[RedOrange] (6.7,1.8) rectangle (8,2);
\fill[RedOrange] (17,1.8) rectangle (18.5,2);
%field oxides:
\fill[DarkGray] (0,2) rectangle (1,4);
\fill[DarkGray] (8.5,2) rectangle (11.5,4);
\fill[DarkGray] (19,2) rectangle (20,4);
% channel stop
\fill[RedOrange] (0,1.5) rectangle (1,2);
\fill[RedOrange] (8.5,1.5) rectangle (11.5,2);
\fill[RedOrange] (19,1.5) rectangle (20,2);
	\end{tikzpicture}
\end{center}

\subsubsection{Drive-in}
\begin{center}
	\begin{tikzpicture}[node distance = 3cm, auto, thick,scale=\CrossSectionOnly, every node/.style={transform shape}]
		% substrate
\fill[YellowOrange] (0,0) rectangle (20,2);
\node at (2,0.5) {Si (p-type)};
% n-well
\fill[Goldenrod] (1.25,0.75) rectangle (8.25,2);
\node at (4.75,1) {N-Well};
% gate oxide
\fill[LightGray] (4.8,2) rectangle (6.7,2.1);
\fill[LightGray] (13.3,2) rectangle (15.2,2.1);
% gate poly
\fill[BrickRed] (4.8,2.1) rectangle (6.7,2.2);
\fill[BrickRed] (13.3,2.1) rectangle (15.2,2.2);
% dioxide
\fill[NormalGray] (0,2) rectangle (3.5,2.7);
\fill[NormalGray] (8,2) rectangle (13.3,2.7);
\fill[NormalGray] (13.3,2.2) rectangle (15.2,2.7);
\fill[NormalGray] (15.2,2) rectangle (17,2.7);
\fill[NormalGray] (18.5,2) rectangle (20,2.7);
% dioxide 2nd
\fill[NormalGray] (0,2.8) rectangle (3.5,3.3);
\fill[NormalGray] (3.5,2) rectangle (4.8,2.8);
\fill[NormalGray] (4.8,2.3) rectangle (6.7,2.8);
\fill[NormalGray] (6.7,2) rectangle (8,2.8);
\fill[NormalGray] (8,2.8) rectangle (8.5,3.3);
\fill[NormalGray] (11.5,2.8) rectangle (17,3.3);
\fill[NormalGray] (17,2) rectangle (18.5,2.8);
\fill[NormalGray] (18.5,2.8) rectangle (19,3.3);

\fill[DarkGray] (0,4.1) rectangle (1,4.2);
\fill[DarkGray] (8.5,4.1) rectangle (11.5,4.2);
\fill[DarkGray] (19,4.1) rectangle (20,4.2);
% p+
\fill[RedOrange] (0,4) rectangle (1,4.1);
\fill[RedOrange] (1,2.7) rectangle (3.5,2.8);
\fill[RedOrange] (8,2.7) rectangle (8.5,2.8);
\fill[RedOrange] (8.5,4) rectangle (11.5,4.1);
\fill[RedOrange] (11.5,2.7) rectangle (17,2.8);
\fill[RedOrange] (18.5,2.7) rectangle (19,2.8);
\fill[RedOrange] (19,4) rectangle (20,4.1);
\fill[RedOrange] (3.5,1.8) rectangle (4.8,2);
\fill[RedOrange] (4.8,2.2) rectangle (6.7,2.3);
\fill[RedOrange] (6.7,1.8) rectangle (8,2);
\fill[RedOrange] (17,1.8) rectangle (18.5,2);
%field oxides:
\fill[DarkGray] (0,2) rectangle (1,4);
\fill[DarkGray] (8.5,2) rectangle (11.5,4);
\fill[DarkGray] (19,2) rectangle (20,4);
% channel stop
\fill[RedOrange] (0,1.5) rectangle (1,2);
\fill[RedOrange] (8.5,1.5) rectangle (11.5,2);
\fill[RedOrange] (19,1.5) rectangle (20,2);
	\end{tikzpicture} \\
	\includegraphics[scale=0.01]{down_arrow.png} \\
	\begin{tikzpicture}[node distance = 3cm, auto, thick,scale=\CrossSectionOnly, every node/.style={transform shape}]
		% substrate
\fill[YellowOrange] (0,0) rectangle (20,2);
\node at (2,0.5) {Si (p-type)};
% n-well
\fill[Goldenrod] (1.25,0.75) rectangle (8.25,2);
\node at (4.75,1) {N-Well};
% gate oxide
\fill[LightGray] (4.8,2) rectangle (6.7,2.1);
\fill[LightGray] (13.3,2) rectangle (15.2,2.1);
% gate poly
\fill[BrickRed] (4.8,2.1) rectangle (6.7,2.2);
\fill[BrickRed] (13.3,2.1) rectangle (15.2,2.2);
% dioxide
\fill[NormalGray] (0,2) rectangle (3.5,2.7);
\fill[NormalGray] (8,2) rectangle (13.3,2.7);
\fill[NormalGray] (13.3,2.2) rectangle (15.2,2.7);
\fill[NormalGray] (15.2,2) rectangle (17,2.7);
\fill[NormalGray] (18.5,2) rectangle (20,2.7);
% dioxide 2nd
\fill[NormalGray] (0,2.8) rectangle (3.5,3.3);
\fill[NormalGray] (3.5,2) rectangle (4.8,2.8);
\fill[NormalGray] (4.8,2.3) rectangle (6.7,2.8);
\fill[NormalGray] (6.7,2) rectangle (8,2.8);
\fill[NormalGray] (8,2.8) rectangle (8.5,3.3);
\fill[NormalGray] (11.5,2.8) rectangle (17,3.3);
\fill[NormalGray] (17,2) rectangle (18.5,2.8);
\fill[NormalGray] (18.5,2.8) rectangle (19,3.3);

\fill[DarkGray] (0,4.1) rectangle (1,4.2);
\fill[DarkGray] (8.5,4.1) rectangle (11.5,4.2);
\fill[DarkGray] (19,4.1) rectangle (20,4.2);
% p+
\fill[RedOrange] (0,4) rectangle (1,4.1);
\fill[RedOrange] (1,2.7) rectangle (3.5,2.8);
\fill[RedOrange] (8,2.7) rectangle (8.5,2.8);
\fill[RedOrange] (8.5,4) rectangle (11.5,4.1);
\fill[RedOrange] (11.5,2.7) rectangle (17,2.8);
\fill[RedOrange] (18.5,2.7) rectangle (19,2.8);
\fill[RedOrange] (19,4) rectangle (20,4.1);
\fill[RedOrange] (3.4,1.2) rectangle (4.9,2);
\fill[RedOrange] (4.8,2.2) rectangle (6.7,2.3);
\fill[RedOrange] (6.6,1.2) rectangle (8.1,2);
\fill[RedOrange] (16.9,1.2) rectangle (18.7,2);
%field oxides:
\fill[DarkGray] (0,2) rectangle (1,4);
\fill[DarkGray] (8.5,2) rectangle (11.5,4);
\fill[DarkGray] (19,2) rectangle (20,4);
% channel stop
\fill[RedOrange] (0,1.5) rectangle (1,2);
\fill[RedOrange] (8.5,1.5) rectangle (11.5,2);
\fill[RedOrange] (19,1.5) rectangle (20,2);
	\end{tikzpicture}
\end{center}

\subsubsection{Oxide removal}
In the final step before the  next mask we need to remove the impure oxide layer for further processing of the wafer.

\begin{center}
	\begin{tikzpicture}[node distance = 3cm, auto, thick,scale=\CrossSectionOnly, every node/.style={transform shape}]
		% substrate
\fill[YellowOrange] (0,0) rectangle (20,2);
\node at (2,0.5) {Si (p-type)};
% n-well
\fill[Goldenrod] (1.25,0.75) rectangle (8.25,2);
\node at (4.75,1) {N-Well};
% gate oxide
\fill[LightGray] (4.8,2) rectangle (6.7,2.1);
\fill[LightGray] (13.3,2) rectangle (15.2,2.1);
% gate poly
\fill[BrickRed] (4.8,2.1) rectangle (6.7,2.2);
\fill[BrickRed] (13.3,2.1) rectangle (15.2,2.2);
% dioxide
\fill[NormalGray] (0,2) rectangle (3.5,2.7);
\fill[NormalGray] (8,2) rectangle (13.3,2.7);
\fill[NormalGray] (13.3,2.2) rectangle (15.2,2.7);
\fill[NormalGray] (15.2,2) rectangle (17,2.7);
\fill[NormalGray] (18.5,2) rectangle (20,2.7);
% dioxide 2nd
\fill[NormalGray] (0,2.8) rectangle (3.5,3.3);
\fill[NormalGray] (3.5,2) rectangle (4.8,2.8);
\fill[NormalGray] (4.8,2.3) rectangle (6.7,2.8);
\fill[NormalGray] (6.7,2) rectangle (8,2.8);
\fill[NormalGray] (8,2.8) rectangle (8.5,3.3);
\fill[NormalGray] (11.5,2.8) rectangle (17,3.3);
\fill[NormalGray] (17,2) rectangle (18.5,2.8);
\fill[NormalGray] (18.5,2.8) rectangle (19,3.3);

\fill[DarkGray] (0,4.1) rectangle (1,4.2);
\fill[DarkGray] (8.5,4.1) rectangle (11.5,4.2);
\fill[DarkGray] (19,4.1) rectangle (20,4.2);
% p+
\fill[RedOrange] (0,4) rectangle (1,4.1);
\fill[RedOrange] (1,2.7) rectangle (3.5,2.8);
\fill[RedOrange] (8,2.7) rectangle (8.5,2.8);
\fill[RedOrange] (8.5,4) rectangle (11.5,4.1);
\fill[RedOrange] (11.5,2.7) rectangle (17,2.8);
\fill[RedOrange] (18.5,2.7) rectangle (19,2.8);
\fill[RedOrange] (19,4) rectangle (20,4.1);
\fill[RedOrange] (3.4,1.2) rectangle (4.9,2);
\fill[RedOrange] (4.8,2.2) rectangle (6.7,2.3);
\fill[RedOrange] (6.6,1.2) rectangle (8.1,2);
\fill[RedOrange] (16.9,1.2) rectangle (18.7,2);
%field oxides:
\fill[DarkGray] (0,2) rectangle (1,4);
\fill[DarkGray] (8.5,2) rectangle (11.5,4);
\fill[DarkGray] (19,2) rectangle (20,4);
% channel stop
\fill[RedOrange] (0,1.5) rectangle (1,2);
\fill[RedOrange] (8.5,1.5) rectangle (11.5,2);
\fill[RedOrange] (19,1.5) rectangle (20,2);
	\end{tikzpicture} \\
	\includegraphics[scale=0.01]{down_arrow.png} \\
	\begin{tikzpicture}[node distance = 3cm, auto, thick,scale=\CrossSectionOnly, every node/.style={transform shape}]
		% substrate
\fill[YellowOrange] (0,0) rectangle (20,2);
% n-well
\fill[Goldenrod] (1.25,0.75) rectangle (8.25,2);
% source
\fill[RedOrange] (0,1.5) rectangle (1,2);
\fill[RedOrange] (3.5,1) rectangle (5,2);
\fill[RedOrange] (6.5,1) rectangle (8,2);
\fill[RedOrange] (8.5,1.5) rectangle (11.5,2);
\fill[RedOrange] (17,1) rectangle (18.5,2);
\fill[RedOrange] (19,1.5) rectangle (20,2);
% gate oxide
\fill[LightGray] (4.8,2) rectangle (6.7,2.1);
\fill[LightGray] (13.3,2) rectangle (15.2,2.1);
% gate poly
\fill[BrickRed] (4.8,2.1) rectangle (6.7,2.2);
\fill[BrickRed] (13.3,2.1) rectangle (15.2,2.2);
%field oxides:
\fill[DarkGray] (0,2) rectangle (1,4);
\fill[DarkGray] (8.5,2) rectangle (11.5,4);
\fill[DarkGray] (19,2) rectangle (20,4);
% channel stop
\fill[RedOrange] (0,1.5) rectangle (1,2);
\fill[RedOrange] (8.5,1.5) rectangle (11.5,2);
\fill[RedOrange] (19,1.5) rectangle (20,2);
	\end{tikzpicture}
\end{center}

For this purpose we use a wet etching method selectively etching away the before grown, each approximately 250nm thick layers of oxide, together 500nm.
Because we don't want to reduce the thick oxide of the field oxide area too much we limit the etching time  here, so that we at worst etch maybe a 10% of it.
Since we have grown it to be around 750nm we can afford to loose around 75nm or so.

